% Source: https://artofproblemsolving.com/community/c3h1582862_17_geometry_questions
% Stable identity: aops-topic-1582862-q01
Each tip moves on a circle of radius $26$ feet:

\begin{asy}
usepackage("amsmath");
size(4.5cm);
defaultpen(fontsize(10pt));
pair O=(0,0);real r=2.6;
draw(circle(O,r),dashed+gray);
for(int i=0;i<3;++i) draw(O--r*dir(90+120*i),linewidth(1.6));
dot(O);label("$26$ ft",0.55*r*dir(210),dir(300));
draw(arc(O,r+0.25,-20,20),Arrow(TeXHead));label("$400$ rpm",(r+0.3)*dir(0),E);
label("\mbox{}",(min(currentpicture,user=true)+max(currentpicture,user=true))/2,black); // keeps the text after the figure black when LaTeX typesets the labels inline
\end{asy}

In one revolution a tip travels the circumference $2\pi\cdot26=52\pi\approx163.363$ feet. At $400$ revolutions per minute it travels $400\cdot52\pi$ feet per minute, i.e.
\[
\dfrac{400\cdot52\pi}{60}\approx\dfrac{65{,}345.1}{60}\approx1089.08\text{ feet per second},
\]
so the tip speed is about $\boxed{1089\text{ feet per second}}$.
